\subsubsection{The Sobolev Imbedding Theorem}\label{subsec:sobolev-imbedding-theorem}

We begin with a definition of the term ``imbedding.''

\begin{Definition}
\label{def:continuous-imbedding}
Let $X$ and $Y$ be Banach spaces, we say $X$ is \textbf{continuously imbedded} in $Y$ and write\[X\hookrightarrow Y\]if $X\subset Y$ and there is a constant $C$ such that
\begin{equation}\label{eq:continuous-embedding-norm-bound}
\|u\|_Y\leq C\|u\|_X\quad\forall u\in X.
\end{equation}

\end{Definition}

\begin{Remark}
\label{rem:continuous-imbedding-bounded-identity}
If $X\subset Y$ then the identity operator $I:X\to Y$ is well defined. Condition \eqref{eq:continuous-embedding-norm-bound} says that the output of the operator is bounded by a constant times the norm of the input. Thus, the identity is a bounded operator.
\end{Remark}

The Sobolev imbedding theorem says that if $\Omega$ is nice, then for sufficiently large $k$, $H^k(\Omega)$ is imbedded in a space of bounded, continuous functions. We first consider the case $\Omega=\mathbb{R}^m$. As a preparation, we state the following lemma.

\begin{Lemma}
\label{lem:integrable-fourier-transform-continuity}
If $\mathcal{F}[u]\in L^1(\mathbb{R}^m)$, then $u$ is a continuous bounded function. Moreover, $\|u\|_{\infty}\leq(2\pi)^{-m/2}\|\mathcal{F}[u]\|_1$.
\end{Lemma}

\begin{Proof}
If $\mathcal{F}[u]\in\mathcal{S}(\mathbb{R}^m)$, then $u\in\mathcal{S}(\mathbb{R}^m)$ and
\begin{equation}\label{eq:integrable-fourier-transform-uniform-bound}
\begin{aligned}
\|u\|_{\infty}&=\sup_{\bm{x}\in\mathbb{R}^m}(2\pi)^{-m/2}\left|\int_{\mathbb{R}^m}\exp(i\bm{\xi}\cdot\bm{x})\mathcal{F}[u](\bm{\xi})\,\mathrm{d}\bm{\xi}\right|\\
&\leq(2\pi)^{-m/2}\|\mathcal{F}[u]\|_1.
\end{aligned}
\end{equation}
Since $\mathcal{S}(\mathbb{R}^m)$ is dense in $L^1(\mathbb{R}^m)$, the lemma follows by taking limits.
\end{Proof}

\begin{Theorem}{Sobolev imbedding theorem}
\label{thm:sobolev-full-space-continuous-embedding}
Let $s>m/2$. Then\[H^s(\mathbb{R}^m)\hookrightarrow C_b(\mathbb{R}^m).\]That is\[H^s(\mathbb{R}^m)\subset C_b(\mathbb{R}^m)\]and there is a constant $C$ such that $\|u\|_{\infty}\leq C\|u\|_{s,2}$ for every $u\in H^s(\mathbb{R}^m)$.
\end{Theorem}

\begin{Proof}
Let $u\in H^s(\mathbb{R}^m)$. Then $\mathcal{F}[u]\in L_s^2(\mathbb{R}^m)$. By the previous lemma, it therefore suffices to show that $L_s^2(\mathbb{R}^m)$ is continuously imbedded in $L^1(\mathbb{R}^m)$. Let $w(\bm{\xi})=1+|\bm{\xi}|^{2s}$. Then Hölder\textquotesingle s inequality yields
\begin{equation}\label{eq:sobolev-embedding-fourier-holder-weight}
\|\mathcal{F}[u]\|_1\leq\|w^{1/2}\mathcal{F}[u]\|_2\|w^{-1/2}\|_2.
\end{equation}
Since we have assumed $s>m/2$, $w^{-1/2}$ is in $L^2(\mathbb{R}^m)$.
\end{Proof}

\begin{Corollary}
\label{cor:sobolev-extension-continuity}
Assume that $k>m/2$ and that $\Omega$ has the $k$-extension property. Then\[H^k(\Omega)\hookrightarrow C_b(\overline\Omega).\]If $k>(m/2)+j$, then\[H^k(\Omega)\hookrightarrow C_b^j(\overline\Omega).\]
\end{Corollary}

\begin{Remark}
\label{rem:sobolev-embedding-optimal-threshold}
In Theorem~\ref{thm:sobolev-full-space-continuous-embedding}, the bound $s>m/2$ is optimal; cf. Problem~\ref{ex:sobolev-embedding-critical-counterexample}.
\end{Remark}

\Needspace{4\baselineskip}
\begin{Example}
\label{ex:sobolev-jump-discontinuity}
In Example~\ref{ex:heaviside-absolute-value-sobolev-membership} we noted that the Heaviside function $H(x)$ was not in $H^1(-1,1)$. The same argument could be used to show that no function with a jump discontinuity is in $H^1(-1,1)$. The Sobolev imbedding theorem gives a much stronger result, stating that if $s>1/2$ every function in $H^s(-1,1)$ is bounded and continuous.
\end{Example}

\Needspace{4\baselineskip}
\begin{Example}
\label{ex:sobolev-radial-power-embedding}
Let $B$ be the ball of radius one in $\mathbb{R}^m$. In Example~\ref{ex:radial-power-sobolev-membership} we showed that $|\bm{x}|^\alpha\in W^{k,p}(B)$ if and only if $\alpha>k-m/p$. If we set $p=2$ we note that the condition of the Sobolev imbedding theorem requiring $k>m/p$ is exactly the condition that ensures $\alpha>0$ so that our radial function would be continuous.
\end{Example}

\Needspace{4\baselineskip}
\begin{Example}
\label{ex:sobolev-three-dimensional-critical-order}
In $\mathbb{R}^3$ the critical ``number of derivatives'' is $s=3/2$. Functions in $H^1(\mathbb{R}^3)$ are not necessarily continuous, but all functions in $H^2(\mathbb{R}^3)$ are bounded and continuous.
\end{Example}

\begin{Remark}
\label{rem:sobolev-general-p-embeddings}
More general imbedding theorems for $W^{k,p}$ spaces can be established. In particular, it can be shown that\[W^{k,p}(\mathbb{R}^m)\hookrightarrow L^{mp/(m-kp)}(\mathbb{R}^m)\quad\text{for }kp<m,\]and\[W^{k,p}(\mathbb{R}^m)\hookrightarrow C_b(\mathbb{R}^m)\quad\text{for }kp>m,\]and that each of these imbeddings is continuous. Again this extends to more general domains $\Omega$, e.g., those which have an extension property.
\end{Remark}
